% GENERATED FILE. Edit canonical lessons, then run npm run generate:books.
% canonical-source-hash: fnv1a64:9e80b669eda6d5c8
% canonical-lessons: RU-C328-lomtik, RU-C328-kusochek, RU-C328-narezat, RU-C328-natirat, RU-C328-tyorka

\chapter{Slicing and grating}
\label{ch:ru-slicing-and-grating}
\hlchaptermodality{328}

\begin{chapteropening}
\textbf{By the end of this chapter:} \emph{I can describe getting food ready: a slice, a small piece, slicing, grating and a grater.}

Complete the last lesson of chapter 328: I can describe getting food ready: a slice, a small piece, slicing, grating and a grater.
\end{chapteropening}

\section[\texorpdfstring{\ru{ломтик} (lómtik)}{ломтик (lómtik)}]{\ru{ломтик} (lómtik) — a slice}

\label{lesson:RU-C328-lomtik}



\begin{quote}
Before the new one: say the Russian for a treasure, then the Russian for a buried treasure.
\end{quote}

\subsection*{You'll want to know: \ru{ломтик}}
\textbf{\ru{ломтик}} (\emph{lómtik}) — \textquotedblleft{}a slice\textquotedblright{}.

\begin{quote}
\ru{Она} \ru{положила} \ru{на} \ru{тарелку} \ru{два} \ru{ломтика} \ru{сыра}.
\end{quote}

\begin{quote}
\emph{Oná polozhíla na tarélku dva lómtika sýra.}
\end{quote}

\begin{quote}
She put two slices of cheese on the plate.
\end{quote}

Diminutive of \emph{\ru{ломоть}} \textquotedblleft{}a hunk, a slice\textquotedblright{}.

\begin{cousinweb}
\emph{\ru{ломтик}} is a diminutive of \emph{\ru{ломоть}} \textquotedblleft{}a hunk of bread\textquotedblright{}, built on the root of \emph{\ru{ломать}} \textquotedblleft{}to break\textquotedblright{}: a piece broken off. EtymDB lists no parent.
\end{cousinweb}

\begin{grammarlens}[title={What you've built}]
One more word for telling, explaining and describing at B1.
\end{grammarlens}

\subsection*{Your turn}
\begin{itemize}
\raggedright
  \item \emph{Say it:} \emph{lómtik}
  \item \emph{Say it:} \emph{lómtik}, once more
  \item \emph{Say it:} a sentence of your own with \emph{lómtik}
  \item \emph{Recall:} say the Russian for a treasure, then the Russian for a buried treasure, then say \emph{lómtik} again
\end{itemize}

\begin{tcolorbox}[breakable,colback=teal!4,colframe=teal!35!black,title={Before you move on}]
What does \ru{ломтик} mean? (\textquotedblleft{}A slice\textquotedblright{}.) Say it once more.
\end{tcolorbox}
\section[\texorpdfstring{\ru{кусочек} (kusóchek)}{кусочек (kusóchek)}]{\ru{кусочек} (kusóchek) — a little piece, a bit}

\label{lesson:RU-C328-kusochek}



\begin{quote}
Before the new one: say the Russian for a buried treasure, then the Russian for a slice.
\end{quote}

\subsection*{You'll want to know: \ru{кусочек}}
\textbf{\ru{кусочек}} (\emph{kusóchek}) — \textquotedblleft{}a little piece, a bit\textquotedblright{}.

\begin{quote}
\ru{Он} \ru{съел} \ru{кусочек} \ru{торта}.
\end{quote}

\begin{quote}
\emph{On syel kusóchek tórta.}
\end{quote}

\begin{quote}
He ate a piece of cake.
\end{quote}

Diminutive of \emph{\ru{кусок}} \textquotedblleft{}a piece\textquotedblright{}; genitive \emph{\ru{кусочка}}.

\begin{cousinweb}
\emph{\ru{кусочек}} is a diminutive of \emph{\ru{кусок}} \textquotedblleft{}a piece\textquotedblright{}, built on \emph{\ru{кус}} \textquotedblleft{}a mouthful\textquotedblright{}, which is inherited from Proto-Slavic \emph{kǫsŭ} \textquotedblleft{}piece\textquotedblright{}, from \emph{kǫsati} \textquotedblleft{}to bite\textquotedblright{} (EtymDB), the root of \emph{\ru{кусать}}.
\end{cousinweb}

\begin{grammarlens}[title={What you've built}]
One more word for telling, explaining and describing at B1.
\end{grammarlens}

\subsection*{Your turn}
\begin{itemize}
\raggedright
  \item \emph{Say it:} \emph{kusóchek}
  \item \emph{Say it:} \emph{kusóchek}, once more
  \item \emph{Say it:} a sentence of your own with \emph{kusóchek}
  \item \emph{Recall:} say the Russian for a buried treasure, then the Russian for a slice, then say \emph{kusóchek} again
\end{itemize}

\begin{tcolorbox}[breakable,colback=teal!4,colframe=teal!35!black,title={Before you move on}]
What does \ru{кусочек} mean? (\textquotedblleft{}A little piece, a bit\textquotedblright{}.) Say it once more.
\end{tcolorbox}
\section[\texorpdfstring{\ru{нарезать} (narézat)}{нарезать (narézat)}]{\ru{нарезать} (narézat) — to slice, to cut up}

\label{lesson:RU-C328-narezat}



\begin{quote}
Before the new one: say the Russian for a slice, then the Russian for a little piece.
\end{quote}

\subsection*{You'll want to know: \ru{нарезать}}
\textbf{\ru{нарезать}} (\emph{narézat}) — \textquotedblleft{}to slice, to cut up\textquotedblright{}.

\begin{quote}
\ru{Мама} \ru{нарезала} \ru{хлеб}.
\end{quote}

\begin{quote}
\emph{Máma narézala khleb.}
\end{quote}

\begin{quote}
Mum sliced the bread.
\end{quote}

Perfective, stressed on \emph{-\ru{ре}-}: \emph{\ru{нарежу}}. The imperfective is \emph{\ru{нарезать}} (\emph{narezát}), stressed on the ending.

\begin{cousinweb}
\emph{\ru{нарезать}} is built on \emph{\ru{на}-} and \emph{\ru{резать}} \textquotedblleft{}to cut\textquotedblright{}, inherited from Proto-Slavic \emph{rězati} (EtymDB).
\end{cousinweb}

\begin{grammarlens}[title={What you've built}]
One more word for telling, explaining and describing at B1.
\end{grammarlens}

\subsection*{Your turn}
\begin{itemize}
\raggedright
  \item \emph{Say it:} \emph{narézat}
  \item \emph{Say it:} \emph{narézat}, once more
  \item \emph{Say it:} a sentence of your own with \emph{narézat}
  \item \emph{Recall:} say the Russian for a slice, then the Russian for a little piece, then say \emph{narézat} again
\end{itemize}

\begin{tcolorbox}[breakable,colback=teal!4,colframe=teal!35!black,title={Before you move on}]
What does \ru{нарезать} mean? (\textquotedblleft{}To slice, to cut up\textquotedblright{}.) Say it once more.
\end{tcolorbox}
\section[\texorpdfstring{\ru{натирать} (natirát)}{натирать (natirát)}]{\ru{натирать} (natirát) — to grate; to rub}

\label{lesson:RU-C328-natirat}



\begin{quote}
Before the new one: say the Russian for a little piece, then the Russian for to slice.
\end{quote}

\subsection*{You'll want to know: \ru{натирать}}
\textbf{\ru{натирать}} (\emph{natirát}) — \textquotedblleft{}to grate; to rub\textquotedblright{}.

\begin{quote}
\ru{Я} \ru{натёр} \ru{сыр} \ru{для} \ru{пиццы}.
\end{quote}

\begin{quote}
\emph{Ya natyór syr dlya pítstsy.}
\end{quote}

\begin{quote}
I grated cheese for the pizza.
\end{quote}

Imperfective; the perfective is \emph{\ru{натереть}} (\emph{naterét}), \emph{\ru{натру}}, past \emph{\ru{натёр}} (\emph{natyór}).

\begin{cousinweb}
\emph{\ru{натирать}} is built on \emph{\ru{на}-} and \emph{\ru{тереть}} \textquotedblleft{}to rub\textquotedblright{}, inherited from Proto-Slavic \emph{terti}, from PIE \emph{terh\textsubscript{1}-} \textquotedblleft{}to rub\textquotedblright{} (EtymDB).
\end{cousinweb}

\begin{grammarlens}[title={What you've built}]
One more word for telling, explaining and describing at B1.
\end{grammarlens}

\subsection*{Your turn}
\begin{itemize}
\raggedright
  \item \emph{Say it:} \emph{natirát}
  \item \emph{Say it:} \emph{natirát}, once more
  \item \emph{Say it:} a sentence of your own with \emph{natirát}
  \item \emph{Recall:} say the Russian for a little piece, then the Russian for to slice, then say \emph{natirát} again
\end{itemize}

\begin{tcolorbox}[breakable,colback=teal!4,colframe=teal!35!black,title={Before you move on}]
What does \ru{натирать} mean? (\textquotedblleft{}To grate; to rub\textquotedblright{}.) Say it once more.
\end{tcolorbox}
\section[\texorpdfstring{\ru{тёрка} (tyórka)}{тёрка (tyórka)}]{\ru{тёрка} (tyórka) — a grater}

\label{lesson:RU-C328-tyorka}



\begin{quote}
Before the new one: say the Russian for to slice, then the Russian for to grate.
\end{quote}

\subsection*{You'll want to know: \ru{тёрка}}
\textbf{\ru{тёрка}} (\emph{tyórka}) — \textquotedblleft{}a grater\textquotedblright{}.

\begin{quote}
\ru{Тёрка} \ru{лежит} \ru{рядом} \ru{с} \ru{раковиной}.
\end{quote}

\begin{quote}
\emph{Tyórka lezhít ryádom s rákovinoy.}
\end{quote}

\begin{quote}
The grater is next to the sink.
\end{quote}

Feminine; genitive plural \emph{\ru{тёрок}} (\emph{tyórok}).

\begin{cousinweb}
\emph{\ru{тёрка}} is built on \emph{\ru{тереть}} \textquotedblleft{}to rub\textquotedblright{}, inherited from Proto-Slavic \emph{terti}, PIE \emph{terh\textsubscript{1}-} (EtymDB), with the ending \emph{-\ru{ка}} for a tool.
\end{cousinweb}

\begin{grammarlens}[title={What you've built}]
One more word for telling, explaining and describing at B1.
\end{grammarlens}

\subsection*{Your turn}
\begin{itemize}
\raggedright
  \item \emph{Say it:} \emph{tyórka}
  \item \emph{Say it:} \emph{tyórka}, once more
  \item \emph{Say it:} a sentence of your own with \emph{tyórka}
  \item \emph{Recall:} say the Russian for to slice, then the Russian for to grate, then say \emph{tyórka} again
\end{itemize}

\begin{tcolorbox}[breakable,colback=teal!4,colframe=teal!35!black,title={Before you move on}]
What does \ru{тёрка} mean? (\textquotedblleft{}A grater\textquotedblright{}.) Say it once more.
\end{tcolorbox}
